\documentclass[12pt,a4paper]{article}
\usepackage[top=22mm,bottom=24mm,left=27mm,right=27mm]{geometry}
\usepackage{amsmath,amssymb,xcolor}
\definecolor{link}{HTML}{205973}
\usepackage[hypertexnames=false,colorlinks=true,linkcolor=link,urlcolor=link]{hyperref}
\hypersetup{pdftitle={The intercept in the idealized sampled program}}
\setlength{\parindent}{0pt}
\setlength{\parskip}{5pt}
\raggedbottom
\newcommand{\R}{\mathbb R}
\newcommand{\lead}[1]{\par\smallskip\noindent\textbf{#1}\enspace}
\begin{document}
{\LARGE\bfseries The intercept bound in program (18)\par}
\medskip
Section 3.2 of Bhat, Farias, Moallemi and Zheng gives a bound on the
intercept of (18). In a two-state model the printed bound can reduce the
optimal value after a single draw. The central order statistics of the sampled Bellman
errors determine exactly which intercepts maximize the objective.

\section*{The sampled problem and its optimal intercepts}
Let \(x_1,\ldots,x_N\), \(N\ge1\), be the sampled states in a finite Markov
decision process with stage costs \(g_{x,a}\), transition probabilities
\(p(x,y,a)\), and discount \(0<\alpha<1\). Each draw contributes one
term, including repeated states. For a real Hilbert space feature map
\(\Phi\), write \(V_z(x)=\langle\Phi(x),z\rangle\) and \(J=V_z+b\).
With an initial distribution \(\nu\) and radii \(C,B\ge0\), the program is
\begin{equation}
\begin{aligned}
 \operatorname{maximize}_{z,b,s}\quad &\nu^\top J-\frac{2}{(1-\alpha)N}\sum_{i=1}^N s_i\\
 \text{subject to}\quad &J(x_i)\le (TJ)(x_i)+s_i,\quad s_i\ge0,\quad
 \|z\|\le C,\quad |b|\le B,
\end{aligned}\tag{18}
\end{equation}
where \((TV)(x)=\min_a\{g_{x,a}+\alpha\sum_y p(x,y,a)V(y)\}\).
The free problem allows any real \(b\).

\lead{Fixed feature weight.}
For fixed \(z\), put \(t_i=((TV_z)(x_i)-V_z(x_i))/(1-\alpha)\).
Since each transition row sums to one, \(T(V_z+b)=TV_z+\alpha b\).
The least feasible slack is therefore \(s_i=(1-\alpha)(b-t_i)_+\),
where \(u_+=\max(u,0)\). The slack penalty is strictly positive,
so every optimum uses these slacks. Substitution, using
\(2u_+=u+|u|\) and \(\nu^\top\mathbf 1=1\), gives the objective
\[
 F(z,b)=\nu^\top V_z+\frac1N\sum_{i=1}^Nt_i
                   -\frac1N\sum_{i=1}^N|b-t_i|.
\]
Sort the breakpoints as \(t_{(1)}\le\cdots\le t_{(N)}\). The complete
set of maximizing free intercepts at this \(z\) is
\[
 I_z=[\,t_{(\lceil N/2\rceil)},\ t_{(\lfloor N/2\rfloor+1)}\,].
\]
Indeed, pair the smallest breakpoint with the largest, the second
with the second largest, and so on. For \(u\le v\),
\(|b-u|+|b-v|\ge v-u\), with equality precisely for \(u\le b\le v\).
Intersect the paired intervals, including the middle singleton for odd
\(N\). Their intersection is exactly \(I_z\), also when breakpoints repeat.

Feasibility alone places no upper limit on a free intercept. Every real
\(b\) is feasible with the slacks displayed above. The bound on the next
page concerns optimizing triples.

\newpage
\section*{A uniform bound for every optimizer}
Put \(K=\max_x\|\Phi(x)\|\) and \(G=\max_{x,a}|g_{x,a}|\). For any
\(\|z\|\le C\), both \(|V_z(x)|\) and \(|\sum_y p(x,y,a)V_z(y)|\)
are at most \(CK\). Thus every feasible action has
\[
 \left|g_{x,a}+\alpha\sum_y p(x,y,a)V_z(y)-V_z(x)\right|
 \le G+(1+\alpha)CK .
\]
Taking the action minimum keeps each breakpoint within \([-R,R]\), where
\[
 \boxed{\displaystyle R=\frac{G+(1+\alpha)CK}{1-\alpha}.}
\]
Every free optimizer has \(b\in I_z\subseteq[-R,R]\). Conversely each
feasible free point is dominated by choosing a median at the same \(z\)
and using its least feasible slacks. Hence \(B\ge R\) preserves the
optimal value and all optimizing triples. Existence follows by projecting
\(z\) onto the span of the finitely many feature vectors. This preserves
the fits and the norm bound. With \(B=R\), the reduced objective is
continuous on a compact parameter set.
If \(B>R\), every maximizing intercept is strictly inside its box.

For an individual \(z\), preserving its best value only requires
\([-B,B]\cap I_z\ne\varnothing\). Equality of the full maximizing
intercept sets requires \(I_z\subseteq[-B,B]\). For example, take two
absorbing states, \(\alpha=1/2\), costs \(-1/2,+1/2\), and zero features.
With one draw of each state, \(I_z=[-1,1]\). A box of radius \(1/2\)
keeps the best value and removes some optimal intercepts.

\section*{A counterexample to the printed threshold}
Let \(\alpha=1/2\), \(C=1\), and \(\Phi(u)=-1,\ \Phi(v)=1\).
There is one action, zero cost, and both states move to \(v\). Start
uniformly and take \(N=1\). The discounted occupancy assigns mass
\((1-\alpha)\nu(u)=1/4\) to \(u\), since only time zero can visit it.
On drawing \(u\), the true objective and constraints reduce to
\[
 \max\ b-4s,\qquad |z|\le1,\qquad s\ge0,\qquad b-3z\le2s.
\]
The free value is \(3\), attained at \((z,b,s)=(1,3,0)\), since \(b-4s\le3z-2s\le3\). The printed threshold
\(2CK+G/(1-\alpha)=2\) permits \(B=5/2\). The boxed value is
\(5/2\), attained at \((1,5/2,0)\), since \(b-4s\le B\).
Here \(R=3\), so equality \(B=R\) can put an optimum on the boundary.
More generally, giving both states cost \(G\), features \(-K,+K\),
and taking \(z=C\) yields breakpoint \(R\) and attains the free
maximum there. This proves sharpness using only \(G,C,K,\alpha\).

{\small\textbf{Sources.}
\href{https://moallemi.com/ciamac/papers/kernel-adp-2012.pdf\#page=12}{Bhat, Farias, Moallemi and Zheng, (18) and \S3.2, pp. 12--13}.
\href{https://github.com/nikhilpb/kadp/blob/main/lean/Offset.lean}{Lean proofs and the two-state witness}.}
\end{document}
